\section{The Role of the Cabinet}

\subsection{Composition of the Cabinet}

The revolutionaries are made up of a coalition of revolutionary organisations, anti-authoritarian and left-wing resistance movements operating within the Southern Cone. As there is no single government or leader of the revolutionaries, the cabinet functions as a space for revolutionaries to coordinate with each other. Each organisation will have their own political networks, capacity to mobilise supporters, and ability to operate covertly under conditions of repression. Although united by opposition to authoritarian rule, these organisations have separate leadership structures, national priorities, and ideological traditions.

Since each revolutionary organisation has and may develop further networks, the cabinet’s practical field of action extends beyond any single country. Even though its centre remains in the Southern Cone, contacts with exile communities, sympathetic organisations, and external actors may widen its political and operational reach.

\subsection{Objectives}

The revolutionaries represent those that seek to resist authoritarian rule and support revolutionary movements from being dismantled by a coordinate repressive system. These revolutionary organizations combine political, military, and clandestine methods in pursuit of their objectives. Although they differ considerably in ideology, structure, and long-term goals, many seek to preserve secure communication networks, protect their members from state repression, expand public support, and undermine the legitimacy and authority of military governments. Some also maintain contacts with foreign governments, exile communities, or other revolutionary movements in order to obtain resources, training, intelligence, or logistical assistance.

The degree of cooperation between revolutionaries varies significantly. Shared opposition to authoritarian regimes encourages collaboration in some cases, while ideological disagreements, national priorities, and competing leadership structures limit it in others. As a result, each organization adopts strategies that reflect its political objectives, operational capabilities, available resources, popular support, and willingness to cooperate with other actors.

\subsubsection{Political Organisation and Mobilisation}

Revolutionary organisations depend on continued participation from workers, students, rural communities, political activists, and other social groups. This means that political education, recruitment, propaganda, labour organisation, and public mobilisation are not secondary to armed activity. These actions may enable revolutionary movements to expand their influence, develop support outside their immediate membership, and transform isolated acts of resistance into a broader political challenge. Without a durable social base, the revolutionaries may struggle to conduct their operations as they will lack the political strength required to produce lasting change in Latin America.

\subsubsection{Intelligence and Counterintelligence}

Intelligence is the process of gathering information, whereas counterintelligence is the process of protecting your own information. Hence, while intelligence should be used to gather data on security forces, surveillance activities, and planned operations for example, counterintelligence is used to protect their members, leadership, and communications. Operations which include infiltrating state institutions, intercepting communications, and spreading disinformation can aid the revolutionary movements in challenging stronger state forces without having to match their military resources.

\subsubsection{Regional and International Coordination}

It is important for revolutionary organisations to coordinate their plans across the Southern Cone to share intelligence, distribute resources, protect threatened members, and continue operating when one national movement has been weakened.

Outside of the Southern Cone, it is possible to contact exile communities, sympathetic organisations, or foreign actors which may provide political publicity, funding, training, humanitarian assistance, or material support. However, this assistance is not guaranteed, and external actors may attach conditions or pursue their own interests. Therefore, Foreign support may strengthen the revolutionaries while also limiting their political or operational independence.

\subsection{Internal Tensions}

Although the revolutionaries are united by their opposition to authoritarian rule, this does not remove any political or strategic disagreements. Each revolutionary organisation has their own beliefs and strategies which can cause disagreements within the cabinet. Some areas where disagreements may occur include whether operations should include armed struggle or public mobilisation, revolutionary movements should be coordinated under a regional command or retain their national independence, and if negotiations with moderate groups are considered useful or unacceptable. Other disputes may also emerge over leadership, the distribution of resources, and which organisation should have the greatest influence within the cabinet.

Furthermore, foreign involvement may cause even more disputes within them. Some organisations may view Soviet, Cuban, or other external support as necessary for survival, while others may see it as a threat to their independence. Similar disagreements may develop over the use of violence, the treatment of suspected informants, and whether immediate action or long-term survival should be prioritised.

\subsection{Relationships to Other Cabinets}

\subsubsection{Cabinet A - Regional Council}

The regional council is the revolutionaries main opponent. The governments represented in this cabinet control military forces, police units, intelligence services, detention systems, and national territory. However, not all representatives hold a dictatorship, some rule within a democratic system, which could pose opportunities with higher ranks.

Even though the regimes share an opposition to revolutionary movements, they are not completely united either. Each government has its own national interests, intelligence priorities, and desire for regional influence. These differences may create mistrust, competition, and weaknesses in the coordination between their security services.

\subsubsection{Cabinet B - CIA}

The CIA is a dangerous but complicated actor for the revolutionaries. The CIA’s main objective is to suppress communism both internationally. However, the CIA may not support every action taken by the participating regimes, especially when these actions can threaten wider American interests or create political risks.

\subsubsection{Cabinet D - KGB}

The KGB could be a possible ally to the revolutionaries, but cooperation is not guaranteed. Although there is a shared opposition to American influence and authoritarian governments, the KGB will also act according to Soviet interests. Some opportunities the KGB could create include political support, intelligence sharing, training or funding.

Seeking aid from the KGB may strengthen the revolutionaries, but it may also reduce their political and operational independence. Different organisations within the revolutionaries may pose challenges over Soviet assistance.
